% !TeX root = ../../Book4.tex
% 纲要标题：### C.1 投射有限群与连续离散模
% 纲要位置：工作记录/计划纲要.md，第 3902 行。
% 纲要细目（写作范围）：
% #### C.1.1 有限商与开正规子群
% #### C.1.2 离散连续模与连续余链
% #### C.1.3 有限商计算与导出函子


\section{投射有限群与连续离散模}
\label{b4:sec:C01}

有限商决定投射有限群的拓扑，但一个连续余链所经过的有限商会随余链变化。连续上同调因此需要把所有有限阶段及其过渡映射一起考虑。本节从\cref{b4:def:continuous-group-cochains}的系数范畴出发，将连续上同调识别为不变量函子的右导出，并说明\cref{b4:prop:continuous-finite-quotients}的有限商计算。

\subsection{有限商与开正规子群}
\label{b4:subsec:C01:01}

投射有限群 \(G\) 是有限离散群的投射极限，取有限群直积的子空间拓扑。群运算逐坐标进行，故连续；相容条件定义闭子集，因此 \(G\) 紧且 Hausdorff。有限多个坐标投影核的交是开正规子群，并组成单位元处的邻域基。这里所需的紧性是有限离散空间直积的紧性；同一拓扑事实已在\subsecref{b4:subsec:1104:03}说明。

\ProseParagraphBreak
开子群 \(H\le G\) 的各陪集覆盖紧空间 \(G\)，因而只有有限多个，故 \([G:H]<\infty\)。它也是闭集，因为补集是其他开陪集之并。取它的有限多个共轭子群之交，便得到包含在 \(H\) 中的开正规子群。反过来，若闭子群有有限指数，则其每个陪集的补集是有限个闭陪集之并，故它是开子群。这些判断中不能把“闭”省掉后随意把抽象有限指数子群当作开子群。

\begin{proposition}\label{b4:prop:C-finite-module-quotient}
设 \(G\) 为投射有限群，\(M\) 为有限交换群。\(G\) 对离散 \(M\) 的作用连续，当且仅当存在开正规子群 \(U\triangleleft G\)，使 \(U\) 在 \(M\) 上平凡作用；此时作用通过有限群 \(G/U\)。此外，从 \(G\) 到有限离散群的群同态连续，当且仅当其核为开子群。
\end{proposition}
\begin{proof}
连续作用使每个 \(m\in M\) 的稳定子群开；将有限多个稳定子群相交，得到作用同态的核，这是开正规子群。反过来，经有限离散商的作用显然连续。对于群同态，各个非空纤维都是核的陪集，故所有纤维开与核开等价，也就与到离散群的连续性等价。
\end{proof}

两类基本例子是
\[
 \widehat{\vphantom{Z}\mathbb Z}=\varprojlim_n\mathbb Z/n\mathbb Z,
 \qquad
 \mathbb Z_p=\varprojlim_r\mathbb Z/p^r\mathbb Z.
\]
第一式按正整数的整除关系取极限，第二式固定素数 \(p\)。整数的像在两者中都稠密，\(1\) 是一个拓扑生成元；但 \(\widehat{\vphantom{Z}\mathbb Z}\) 具有任意阶的有限循环商，\(\mathbb Z_p\) 的有限连续商只能是循环 \(p\)-群。后一限制将直接反映在上同调计算中。
\symidx{\(\widehat{\vphantom{Z}\mathbb Z}\)：整数群的投射有限完备化}

\subsection{离散连续模与连续余链}
\label{b4:subsec:C01:02}

设 \(G\) 为投射有限群。一个离散连续 \(G\)-模是带 \(G\) 作用的交换群 \(M\)，其中每个元素的稳定子群都是开的；这正是\cref{b4:def:continuous-group-cochains}中作用映射连续的等价条件。每个开稳定子群又包含一个开正规子群，所以这里的条件是：对每个 \(m\in M\)，存在固定 \(m\) 的开正规子群 \(U_m\)。这个子群可以随 \(m\) 改变，并不要求存在一个开正规子群固定整个 \(M\)。\cref{b4:prop:C-no-common-action-quotient}将给出没有共同有限作用商的例子。对象不必有限。例如，若 \(K^s\) 是域 \(K\) 的可分闭包，则每个代数元被某个有限扩张的 Galois 子群固定，所以 \((K^s,+)\) 和 \((K^s)^\times\) 都是离散连续 \(G_K\)-模，其中 \(G_K=\operatorname{Gal}(K^s/K)\) 取 Krull 拓扑。

\ProseParagraphBreak
记这些模及 \(G\)-同态所成的范畴为 \(\mathcal D_G\)。核、像和商都仍为离散连续模：子模中的元素保留原稳定子群，商中的元素则由任一代表的开稳定子群固定。因此 \(\mathcal D_G\) 是交换范畴，嵌入全部抽象 \(\mathbb ZG\)-模的函子正合。不过，两范畴的内射对象不能不经检查便混为一谈。

\ProseParagraphBreak
连续余链为 \(C^n_{\mathrm{cts}}(G,M)=\operatorname{Map}_{\mathrm{cts}}(G^n,M)\)，次数零取 \(M\)。为方便独立查阅，非齐次微分写为
\begin{align*}
 (\delta f)(g_1,\ldots,g_{n+1})
 &=g_1f(g_2,\ldots,g_{n+1})\\
 &\quad+\sum_{i=1}^n(-1)^i
 f(g_1,\ldots,g_ig_{i+1},\ldots,g_{n+1})\\
 &\quad+(-1)^{n+1}f(g_1,\ldots,g_n).
\end{align*}
零次公式为 \((\delta m)(g)=gm-m\)。这是\eqref{b4:eq:group-cochain-differential}的同一公式，微分平方为零的代数验证继续适用；连续性则来自乘法、作用与有限求和的连续性。

\begin{proposition}\label{b4:prop:C-cochains-exact}
设 \(G\) 为投射有限群，\(\mathcal D_G\) 为离散连续 \(G\)-模及 \(G\)-同态组成的交换范畴。对每个 \(n\ge0\)，函子 \(C^n_{\mathrm{cts}}(G,-):\mathcal D_G\rightarrow\mathbf{Ab}\) 正合。因此连续群上同调对离散连续模的短正合列给出长正合列。
\end{proposition}
\symidx{\(\mathcal D_G\)：离散连续群作用模的范畴}
\begin{proof}
核和单射逐点计算。对满射 \(M\rightarrow N\)，一个连续余链 \(f:G^n\rightarrow N\) 从紧空间映到离散空间，像只有有限多个值。为每个值选一个 \(M\) 中的原像，在它的开闭纤维上取对应常值，得到连续提升。这只是在提升一张连续函数，不要求它等变，也不要求它满足余圈方程；所选原像也不必固定于原先同一个开子群。次数零就是原满射。于是短正合列给出逐次短正合的余链复形，长正合列来自\cref{b4:thm:cohomology-long-exact}的连接态射构造。
\end{proof}

余链可以提升，并不意味着余圈可以提升为余圈。若 \(0\to J\to M\to N\to0\) 中的余圈 \(f\) 提升为余链 \(\widetilde f\)，则 \(\delta\widetilde f\) 在 \(N\) 中的像为 \(\delta f=0\)，因而取值于 \(J\)；它的上同调类就是连接同态记录的提升障碍。次数零的余圈是不变元素，\cref{b4:prop:C-invariant-lift-obstruction}将具体展示一个元素有提升，却没有不变提升的情形。

\subsection{有限商计算与导出函子}
\label{b4:subsec:C01:03}

\cref{b4:prop:continuous-finite-quotients}已经证明，对任意离散连续 \(G\)-模 \(M\)，
\begin{equation}\label{b4:eq:C-finite-quotient-cohomology}
 H^n_{\mathrm{cts}}(G,M)
 \cong\varinjlim_{U\triangleleft_{\mathrm{open}}G}H^n(G/U,M^U).
\end{equation}
指标按开正规子群的反向包含排列：从 \(U\) 到更小的 \(V\) 使用膨胀及系数包含 \(M^U\subseteq M^V\)。证明中的两个有限性步骤是：连续余链像有限，并且它在某个共同开正规子群的陪集上常值。再缩小这个子群使其固定有限像，余链便通过 \((G/U)^n\) 且取值于 \(M^U\)。\cref{b4:prop:filtered-colimits-cohomology}保证取余极限与取上同调相容。

\ProseParagraphBreak
若 \(M\) 有限，先固定一个作用核 \(U_0\)，以后仅取 \(U\subseteq U_0\) 即可，这些指标共尾，且 \(M^U=M\)。但不能只计算 \(H^n(G/U_0,M)\) 后停止：更小子群带来的膨胀映射可能使原有上同调类消失。

\begin{theorem}\label{b4:thm:C-continuous-derived}
设 \(G\) 为投射有限群，\(\mathcal D_G\) 为离散连续 \(G\)-模及 \(G\)-同态组成的交换范畴。则 \(\mathcal D_G\) 有足够多内射对象，并且对每个 \(n\ge0\)，有
\[
 H^n_{\mathrm{cts}}(G,M)\cong R^n(-)^G(M)
 \qquad(M\in\mathcal D_G).
\]
右边的导出函子在离散连续模范畴内计算。
\end{theorem}
\begin{proof}
先从抽象群模范畴取得内射对象，再使其成为离散连续模。对任意抽象 \(\mathbb ZG\)-模 \(B\)，令 \(D(B)\) 为所有被某个开子群固定的元素所成子模。有限交保证加法封闭，开子群的共轭仍开保证 \(G\) 稳定。每个从离散连续模 \(M\) 到 \(B\) 的同态都落入 \(D(B)\)，给出自然双射
\[
 \operatorname{Hom}_{\mathbb ZG}(M,B)
 \cong\operatorname{Hom}_{\mathcal D_G}(M,D(B)).
\]
所以离散化函子 \(D\) 是正合嵌入 \(\mathcal D_G\rightarrow\mathbb ZG\text{-}\mathbf{Mod}\) 的右伴随。抽象模范畴有足够多内射对象；把 \(M\) 嵌入一个抽象内射模 \(J\)，其像落入 \(D(J)\)。\cref{b4:prop:right-adjoint-injectives}说明 \(D(J)\) 在 \(\mathcal D_G\) 中内射，于是得到足够多内射对象。

再检查离散连续内射对象在每个有限商上的高次上同调。令 \(I\) 是 \(\mathcal D_G\) 中的内射对象。对开正规子群 \(U\)，不变量 \(I^U\) 是 \(G/U\)-模。此处使用另一个伴随：沿商映射的膨胀函子与取 \(U\)-不变量满足
\[
 \operatorname{Hom}_{\mathcal D_G}
 (\operatorname{Inf}_{G/U}^G N,I)
 \cong\operatorname{Hom}_{\mathbb Z[G/U]}(N,I^U).
\]
因为左边的同态必须把 \(N\) 映入 \(I^U\)，这个双射由同一张底群映射给出。膨胀不改变底群而正合，故\cref{b4:prop:right-adjoint-injectives}使 \(I^U\) 成为内射 \(G/U\)-模。有限群上同调是抽象群模不变量函子的右导出，因而 \(H^n(G/U,I^U)=0\) 对 \(n>0\) 成立；代入\eqref{b4:eq:C-finite-quotient-cohomology}，得到 \(H^n_{\mathrm{cts}}(G,I)=0\)。第一个伴随建立足够多内射对象，这个伴随则把高次消失化为各有限商上的消失。

由\cref{b4:prop:C-cochains-exact}，连续上同调构成上同调 \(\delta\)-函子，零次为 \(M^G\)。刚证的内射嵌入使其正次可消去，故\cref{b4:thm:derived-universal}将它自然识别为不变量函子的右导出函子。
\end{proof}

这给出了\subsecref{b4:subsec:1104:03}中使用的导出识别。上述范畴与有限商证明可对照 Weibel\cite[Definition 6.11.8; Lemma 6.11.10; Theorem 6.11.13, pp. 210--213]{Weibel1994HomologicalAlgebra}。如果系数带有 \(p\)-进拓扑，例如把 \(\mathbb Z_p\) 本身作为拓扑模，就已离开 \(\mathcal D_G\)；相关连续上同调还要处理逆极限与系数拓扑，本附录不作无说明的替换。

\begin{proposition}\label{b4:prop:C-no-common-action-quotient}
设 \(p\) 为素数，\(G=\mathbb Z_p\)。对 \(r\ge1\)，令 \(M_r=\mathbb F_p[\mathbb Z/p^r\mathbb Z]\)，使 \(G\) 经模 \(p^r\) 商在正则基上作左平移。则离散交换群
\[
 M=\bigoplus_{r\ge1}M_r
\]
是离散连续 \(G\)-模，但其作用不通过任何有限商。
\end{proposition}
\begin{proof}
每个 \(m\in M\) 只有有限多个非零分量。若这些分量的指标都不超过 \(r_0\)，则开正规子群 \(p^{r_0}\mathbb Z_p\) 在这些分量上平凡作用，故固定 \(m\)。零元素被整个 \(G\) 固定，所以每个元素的稳定子群都开，作用连续。

在 \(M_r\) 上，非零的模 \(p^r\) 剩余类会移动正则基中零剩余类所对应的基向量，故作用核恰为 \(p^r\mathbb Z_p\)。于是整个 \(M\) 上的作用核为
\[
 \bigcap_{r\ge1}p^r\mathbb Z_p=\{0\}.
\]
该等式来自逆极限中的元素由全部模 \(p^r\) 坐标决定。因 \(\mathbb Z_p\) 无限，零子群不可能具有有限指数，也不是开子群。因此不存在一个固定整个 \(M\) 的开正规子群，忠实作用也不可能通过有限群分解。
\end{proof}

\begin{proposition}\label{b4:prop:C-invariant-lift-obstruction}
设 \(p\) 为素数，\(G=\mathbb Z_p\) 经商群 \(C_p\) 作用于 \(M=\mathbb F_p[C_p]\)。令 \(\epsilon:M\to\mathbb F_p\) 为系数和，\(J=\ker\epsilon\)，并给 \(\mathbb F_p\) 平凡作用。则离散连续模的短正合列
\[
 0\longrightarrow J\longrightarrow M
 \xrightarrow{\epsilon}\mathbb F_p\longrightarrow0
\]
在每个连续余链次数仍为短正合列，但 \(M^G\to\mathbb F_p\) 是零映射。连接同态将 \(1\) 送到非零类
\[
 [g\mapsto ge-e]\in H^1_{\mathrm{cts}}(G,J),
\]
其中 \(e\) 是群代数中单位元素所对应的基向量。
\end{proposition}
\begin{proof}
正则平移作用使不变向量的全部系数相等，所以
\[
 M^G=\mathbb F_p\cdot\sum_{h\in C_p}h.
\]
这个基向量和的系数和为 \(p=0\)，故不变量映射为零。连续余链上的短正合性由\cref{b4:prop:C-cochains-exact}给出。

用零次余链 \(e\) 提升 \(1\)，其微分是 \(g\mapsto ge-e\)。因 \(\epsilon(ge-e)=0\)，它取值于 \(J\)；又通过有限商 \(C_p\)，故连续。这正是连接同态的代表。若存在 \(j\in J\) 使 \(ge-e=gj-j\) 对所有 \(g\) 成立，则 \(e-j\) 是系数和为 \(1\) 的不变向量，与刚才的零映射结论矛盾。所以该类非零。提升余链后的不变性误差在此不能通过修改提升消除。
\end{proof}
